
\section{Realisability with fixed points}
\label{sec:realisability}
\todonew{
\begin{itemize}
    \item include proof of $\Pi_2$ conservativity of $\muPA$ over $\muHA$ in main text?
    \item switch typesetting of predicate $\pNat$ and type $\nat$ macros?
    \item could make further remarks about abstract realisability just after definition of judgement (some commented copy there).
    \item check and try improve remark about comparison to system F.
    \item at same time can simplify the typed version of $\typedclomuLJneg$ from earlier to be more intuitive.
    \item for functoriality lemma can omit all noncriticial cases?
    \item add comments about abstract realisability inspired by Berger and Tsuiki's appropach, but point out that our realisability model is different as we do not interpret fixed points as fixed points, but are rather inspired by SO encodings.
\end{itemize}
}


So far we have shown an \emph{upper bound} on the representable functions of $\muLJ$ and $\cmuLJ$, namely that they are all provably total in $\PSCA$. 
In this section we turn our attention to proving the analogous \emph{lower bound}, namely by giving a \emph{realisability} interpretation into $\muLJ$ (in fact typed-$\closure \muLJneg$) from a theory over which $\PSCA$ is conservative.


In particular, we consider a version of first-order arithmetic, $\muPA$, with native fixed point operators, (essentially) introduced by M\"ollerfeld in \cite{Moellerfeld02:phd-thesis}.
Unlike that work, we shall formulate $\muPA$ in a purely first-order fashion in order to facilitate the ultimate realisability interpretation.


% Let us write $\langarith (\vec X)$ for the language of arithmetic augmented by predicate symbols $\vec X$, which we shall refer to as `variables'.
% % 
% $\langmuarith$-formulas are generated just like $\langarith$-formulas with the additional clause:

% \begin{itemize}
%     \item if $\phi$ is a formula with free variables $X,\vec X, x, \vec x$, and in which $X$ occurs positively, and $t $ is a (number) term with free variables $\vec y$ then $t \in \mufo X x \phi$ \oldtodo{notation can be optimised, this is a macro} (sometimes $\mufo X x \phi \, t$) is a formula with free variables $\vec X,\vec x,\vec y$.
% \end{itemize}

% Semantically we construe $\mufo X x \phi$ in the standard model as a bona fide least fixed point of the operator defined by $\phi$. 



\subsection{Language of arithmetic with fixed points}
We consider an extension of $\PA$ whose language is closed under (parameterised) least (and greatest) fixed points.
Throughout this section we shall only use logical symbols among $\land, \limp, \exists , \forall$, without loss of generality.

\begin{convention}
\label{conv:langarith-has-all-pr-symbols}
We shall henceforth assume, without loss of generality, that the language of arithmetic $\langarith$ contains a function symbol for each primitive recursive function definition.
All arithmetic theories we consider, like $\PA$, $\HA$ and their extensions, will contain the defining equational axioms for each of these function symbols.
Note that these formulations are just \emph{definitional extensions} of the usual versions of $\PA$ and $\HA$.
However, they allow us to construe $\Delta_0$ formulas as just equations, simplifying some of the metamathematics herein.
\end{convention}




Let us extend the language of arithmetic $\langarith$ by countably many predicate symbols, written $X,Y$ etc., that we shall refer to as `set variables'. 
In this way we shall identify $\langarith$-formulas with the arithmetical formulas of the language of second-order arithmetic $\langsoarith$, but we shall remain in the first-order setting for self-containment.
As such, when referring to the `free variables' of a formula, we include both set variables and number variables.

$\langmuarith$-formulas are generated just like $\langarith$-formulas with the additional clause:

\begin{itemize}
    \item if $\phi$ is a formula with free variables $X,\vec X, x, \vec x$, and in which $X$ occurs positively, and $t $ is a (number) term with free variables $\vec y$ then $t \in \mufo X x \phi$ \oldtodo{notation can be optimised, this is a macro} (or even $\mu \phi\, t$ when $X,x$ are clear from context) is a formula with free variables $\vec X,\vec x,\vec y$.
\end{itemize}



We also write $t \in \nufo X x \phi(X)$ for $t \notin \mufo X x \lnot \phi(\lnot X)$, a suggestive notation witnessing the De Morgan duality between $\mu$ and $\nu$ (in classical logic).

\begin{remark}
[$\langmuarith$ as a proper extension of $\langsoarith$]
\label{rem:langmuarith-extends-langsoarith}
Again referring to the identification of $\langarith$ formulas with arithmetical $\langsoarith$ formulas, we may construe $\langmuarith$ as a formal extension of $\langsoarith$ by certain relation symbols.
Namely
$\langmuarith$ may be identified with the closure of $\langsoarith$ under:
    \begin{itemize}
        \item if $\phi$ is an \emph{arithmetical} formula with free variables among $\vec X, X, \vec x, x$, and in which $X$ occurs positively, then there is a relation symbol $\mufo X x \phi$  taking $|\vec X|$ set inputs and $|\vec x,x|$ number inputs.
    \end{itemize}
In this case, for arithmetical $\phi(\vec X, X, \vec x, x)$, we simply write, say, $t \in \mufo X x {\phi(\vec A, X, \vec t, x)} $ instead of $\mufo X x {\phi(X,\vec X,x,\vec x)} \vec A \,  \vec t\, t$.
Let us note that this is indeed the route taken by M\"ollerfeld.
Instead we choose to treat the construct $\mu \cdot \lambda \cdot $ syntactically as `binder'.
\end{remark}

Semantically we construe $\mufo X x \phi$ in the intended model as a bona fide least fixed point of the (parametrised) operator defined by $\phi$. 
Namely, in the sense of the remark above, we can set $(\mufo X x \phi)^\nmod (\vec A,\vec a)$ to be the least fixed point of the operator $(\Lambda X \lambda x \phi(\vec A, X, \vec a, x))^\nmod$, for all sets $\vec A$ and numbers $\vec a$.
Note that by simple algebraic reasoning this forces $\nufo X x \phi$ to be interpreted as the analogous \emph{greatest} fixed point in $\nmod$.
\anupam{this is a bit brief, but is fine for now}

% \begin{convention}
% [Suppressing variables]
As in earlier parts of this work, we shall frequently suppress variables in formulas to denote abstractions, e.g.\ for formulas $\phi(X)$ and $\psi(x)$, with $X,x$ clear from context, we may write $\phi(\psi)$ for the formula obtained by replacing each subformula $Xt$ of $\phi(X)$ by $\psi(t)$.
% \end{convention}

\subsection{Theories $\muPA$ and $\muHA$}
Let us expand Peano Arithmetic ($\PA$) to the language $\langmuarith$, i.e.\ by including induction instances for all $\langmuarith$-formulas.
The theory we consider here is equivalent to (the first-order part of) M\"ollerfeld's $\ACA(\langmuarith)$.

\begin{definition}
    [Theory]
    \label{def:muPA}
    The theory $\muPA$ is the extension of $\PA$ by the following axioms for formulas $\phi(X,x)$ and $\psi(x)$:
    \begin{itemize}
        \item $\preaxiom_\phi$: $\forall x (\phi(\mu \phi, y) \limp \mu \phi\,  x)$
        \item $\indaxiom_{\phi,\psi}$: $\forall x (\phi(\psi, x) \limp \psi(x)) \limp \forall x (\mu \phi\, x  \limp \psi(x))$
    \end{itemize}
\end{definition}


We sometimes omit the subscripts of the axiom names above.
We may construe $\muPA$ as a proper fragment of full second-order arithmetic $\SOPA$ by the interpretation:
\begin{equation}
    \label{eq:mufo-as-so-int-pre}
    t \in \mufo X x \phi 
    \quad := \quad
    \forall X (\forall x (\phi(X,x) \limp Xx) \limp Xt)
\end{equation}
The axioms above are readily verified by mimicking a standard textbook algebraic proof of Knaster-Tarski in the logical setting, cf.~Proposition~\ref{prop:kt-for-least-pre-fixed-point}.


M\"ollerfeld's main result was that $\PSCA$ is in fact $\Pin 1 1 $-conservative over his theory $\ACA(\langmuarith)$, and so we have:

\begin{theorem}
    [Implied by \cite{Moellerfeld02:phd-thesis}]
    \label{thm:moellerfeld-psca-cons-over-muPA}
    $\PSCA$ is arithmetically conservative over $\muPA$.
\end{theorem}



\anupam{commented fo-parts below}
% \begin{remark}
%     [On first-order parts]
%     % Note that $\muPA$ is axiomatised in a strictly first-order manner to facilitate our ultimate realisability interpretation, whereas 
%     M\"ollerfeld's theory $\ACA(\langmuarith)$ is defined in the same way as $\muPA$ above, only over $\ACA$ instead of $\PA$, cf.~\ref{rem:langmuarith-extends-langsoarith}.
%     The fact that $\muPA$ exhausts the first-order part of M\"ollerfeld's system, i.e.\ that $\ACA(\langmuarith) $ is arithmetically conservative over $\muPA$, follows from standard techniques.
%     In particular any model of $\muPA$ can be expanded to one of $\ACA(\langmuarith)$ by including just the $\mu$-definable sets.
% \end{remark}


% \subsection{Theory $\muHA$ and a negative translation}
We set $\muHA$ to be the intuitionistic counterpart of $\muPA$.
I.e.\ $\muHA$ is axiomatised by Heyting Arithmetic $\HA$ (extended to the language $\langmuarith$) and the schemes $\preaxiom$ and $\indaxiom$ in Definition~\ref{def:muPA} above.


Once again, we may construe $\muHA$ as a proper fragment of full second-order Heyting Arithmetic $\SOHA$ by \eqref{eq:mufo-as-so-int-pre} above.
By essentially specialising known conservativity results for second-order arithmetic, we may thus show that $\muPA$ and $\muHA$ provably define the same recursive functions on natural numbers.
% Indeed essentially the same result was obtained for M\"ollerfeld's system $\ACA{\langmuarith}$ by Tupailo in \cite{Tupailo04doubleneg}, though we include a self contained argument here for the current setting.
\begin{proposition}
[Implied by \cite{Tupailo04doubleneg}]
\label{thm:muPA-pi02-cons-muHA}
    $\muPA$ (so also $\PSCA$) is $\Pin 0 2$-conservative over $\muHA$.
\end{proposition}


We give a self-contained argument for this result in \Cref{sec:conservativity-via-double-negation}, essentially by composing the Friedman-Dragalin $A$-translation and G\"odel-Gentzen negative translation. 


\anupam{commented proof idea below}
% \begin{proof}
%     [Proof idea of \Cref{thm:muPA-pi02-cons-muHA}]
%     We specialise the $\Pin 02$-conservativity of $\SOPA$ over $\SOHA$, using the Friedman-Dragalin $A$-translation and G\"odel-Gentzen negative translation, to $\langmuarith$ and $\muPA$ and $\muHA$.
%     We give a self-contained proof in Appendix~\ref{sec:conservativity-via-double-negation}.
% \end{proof}



% \subsection{Target of realisability}
% We shall consider a rudimentary programming language, with both typed and untyped versions, that suffices as the target of our realisability interpretation.

% \todo{make this all a theorem of $\closure \muLJneg$}

% \oldtodo{need to calibrate this with Gianluca's part}

% Throughout we shall consider only types formed from $\N, \times, \to$ and a special restriction of the fixed point type $ \muN X $.\oldtodo{this is to avoid complicating the realisability argument. We would need both product types (fine) and a version of realisability commuting with quantifiers, which changes emphasis. Instead devolve work to the proof theory part and interpret $\muN X \, \sigma$ as $\mu X (\N \times \sigma)$.}
% A type context is just a type allowing an additional symbol $\cdot$ (we use same metavariables for types and type contexts).
% A type context is positive or negative if $\cdot$ occurs in (only) positive or negative position, respectively.
% We write $\sigma(\tau) := \sigma[\tau/\cdot]$.


% \anupam{to fix clash between target of realisability and typed $\closure \muLJneg$, can change `type of realiser' to decompose $\forall$ step depending on matrix.}
% \newcommand{\Fmu}{F \hspace{-.35em}\restriction  \hspace{-.3em} \mu}
% \begin{definition}
%     [An untyped calculus]
%     Terms are generated by (untyped) application from constants $K,S,I$ etc.\ and $\pair{\cdot}{\cdot}$ and $\projl$ and $\projr$ and $\iterN$ and $\inN {\sigma}$ for $\sigma$ a positive type context, and positive type contexts $\sigma^+$ and negative type contexts $\sigma^-$.
%     The corresponding theory consists of usual equations for $K,S,I$ etc.\ and:
%     \[
%     \begin{array}{r@{\ = \ }l}
%         \projl \pair {x_0}{x_1} & x_0\\
%         \projr \pair {x_0}{x_1} & x_1\\
%          \iterN \, s \, x & x\, s \\
%          \inN {\sigma} \, a\,  x\, h & ha \, (\sigma  (\iter{}\,  h)\, x) \\
%        \cdot \, t & t \\
%        \N t & I\\
%          X t & I   \\
%          (\sigma_0 \times \sigma_1)\, t\, u  & \pair{\sigma_0 t (\projl u)}{\sigma_1 t (\projr u)}\\
%          (\sigma_0 \to \sigma_1)\,  t\, x\, y & \sigma_1\, t \, (x\, (\sigma_0\, t\, y)) \\
%          (\muN X \, \sigma ) \, t & ((\N \to \sigma \to X)\to X)\,  t 
%     \end{array}
%     \]
%     \oldtodo{
%     (NB: the equations for positive and negative type contexts are the same!)}
%     (We also have constants $0, \succ, \rec{} $ satisfying usual equations)
% \end{definition}

% NB: Clearly this language is Turing complete since it contains $K $ and $S$.

% \begin{definition}
%     [Type assignment]
%     We assign types by:
%     \begin{itemize}
%     \item $\projl : (\sigma_0 \times \sigma_1) \to \sigma_0$
%     \item $\projr: (\sigma_0 \times \sigma_1) \to \sigma_1$
%     \item $\pair \cdot \cdot : \sigma \to \tau \to (\sigma \times \tau)$
%         \item $\iterN : (\N \to \sigma(\tau) \to \tau) \to \muN X\,  \sigma(X) \to \tau$ \oldtodo{could consider allowing $\iterN$ to take a numerical argument after step function, to easily realise universal closure of succedent.}
%         \item $\inN \sigma : \N \to \sigma(\muN X\, \sigma(X)) \to \muN X\, \sigma(X)$
%         \item $\sigma^+ : (\tau \to \tau') \to \sigma(\tau) \to \sigma(\tau')$
%         \item $\sigma^- : (\tau \to \tau') \to \sigma(\tau') \to \sigma(\tau)$
%         \item if $t:\sigma \to \tau$ and $s:\sigma$ then $ts : \tau$.
%         \item if $t:\muN X\, \sigma(X)$ and $s:\N \to \sigma(\tau) \to \tau$ then $ts:\tau$.
%         \item 
%         if $t:(\N \to \sigma \to X) \to X$ (and $\sigma$ is positive in $X$) then $t:\muN X \, \sigma$.\oldtodo{A: I added this after thinking about typing; this is one of the interesting cases for embedding into $\muLJ$. As written, types may be open.}
%     \end{itemize}
% \end{definition}

\anupam{commented `target of realisability' above}



% \subsection{Relativisation to $\Nat$}
% It will be convenient for our realisability argument to work with a notion of `abstract' realisability, where realisability commutes with quantifiers for a fragment of the language where all quantifiers are relativised.
% The reason for this is that, a priori, all quantifiers of arithmetic are relativised to $\Nat$, but construing so for the fixed point axioms leads to type mismatch during realisability. We shall make this clear with some remarks throughout, but in this subsection we carry out such a reduction explicitly.

% To avoid confusion with our earlier type $\nat $ for natural numbers, let us write $\pnat \dfn \mu X\lambda x (x=0 \lor \exists y (Xy \land x=\succ y)$. 
% % The notation is suggestive of the previous use of the type $\nat$, but should cause no confusion.
% We also write $\pnatinv(X,x)$ for the matrix of this fixed point, i.e.\ $x=0 \lor \exists y (Xy \land x=\succ y )$.
% \todonew{switch predicate $\pnat$ and type $\nat$ macros?}


% \begin{definition}
%     For formulas $\phi$ of $\langmuarith$ we define $\nattrans \phi$ by:
%     \begin{itemize}
%         \item $ \nattrans{(s=t)} \dfn s=t $
%         \item $ \nattrans{(t \in\mu X \lambda x \phi)} \dfn t \in \mu X \lambda x (\pnat x \land \nattrans\phi)$
%         \item $\nattrans {(\phi \star \psi)} \dfn \nattrans \phi \star \nattrans \psi $ for $\star \in\{\lor, \land, \limp\}$
%         \item $\nattrans {(\exists x \phi)} \dfn \exists x (\pnat x \land \phi)$
%         \item $\nattrans{(\forall x \phi)}\dfn \forall x (\pnat x \limp \phi)$.
%     \end{itemize}
% \end{definition}


% We extend this translation to (definable) predicates, e.g.\ writing $\nattrans \phi(\vec X,\vec x)\dfn \nattrans{\phi(\vec X,\vec x)}$ and, in particular, $\nattrans{(\mu\phi)}x\dfn \nattrans{(x \in \mu\phi)}$.
% We shall employ standard abbreviates for relativised quantifiers, in particular writing $\exists x \in \pNat \, \phi \dfn \exists x (\pnat x \land \phi)$ and $\forall x \in \pNat \, \phi \dfn \forall x (\pnat x \limp \phi) $. 


% Under this interpretation $\muHA$ can be embedded into its induction-free part, a modest extension of pure intuitionistic logic with fixed points. Let us make this clear.

% Write $\muIFOL$ for the intuitionistic theory given by the basic Robinson axioms of arithmetic \todonew{for simplicity. can also just add all true equations} and the $(\preaxiom)$ and $(\indaxiom)$ axioms from \cref{def:muPA}.
% \todonew{can write, say, $\muHA^-$ instead of $\muIFOL$?}
% Specialising a well-known reduction from second-order arithmetic to pure second-order logic, we have:

% \begin{proposition}
%     If $\muHA \proves \phi$ then $\muIFOL \proves \nattrans \phi$.
% \end{proposition}
% \begin{proof}
%     The Robinson axioms are already part of $\muIFOL$ and, since they are universal statements, so too their relativisations.
    
%     For each number induction axiom,
%     \begin{equation}
%         \label{eq:n-ind-instance-for-relativisation}
%         \psi(0) \limp \forall x (\psi(x) \limp \psi(\succ x)) \limp \forall x \psi(x)
%     \end{equation}
%     we show that its $\pNat$-relativisation is equivalent to the $(\indaxiom)$ axiom for $\pNat$ with invariant $\pnat x \land \nattrans\psi(x)$. 
%     Namely, reasoning inside $\muIFOL$, we have,
%     \[
%     \begin{array}{rll}
%          & \forall y ( \pnatinv (\pnat \cdot \land \nattrans\psi(\cdot),y) \limp (\pnat y \land \nattrans\psi(y))) \limp \forall x \in \pNat (\pnat x \land \nattrans\psi(x) )
%          & \text{by $(\indaxiom)$ for $\pNat$} 
%          \\
%          \implies & \pnat 0 \limp \nattrans\psi (0) \limp \forall x \in \pNat (  \nattrans\psi(x) \limp (\pnat {\succ x} \land \nattrans\psi(\succ x)) ) \limp \forall x \in \pNat  \, \nattrans\psi(x)
%          & \text{by pure logic}
%          \\ 
%          \implies & \nattrans\psi(0) \limp \forall x \in \pNat \,  (\nattrans\psi(x) \limp \nattrans\psi(\succ x)) \limp \forall x \in \pNat \,  \nattrans\psi(x)
%          & \text{by $(\preaxiom)$ for $\pNat$}
%     \end{array}
%     \]
%     where the last line is just the $\pNat$-translation of \eqref{eq:n-ind-instance-for-relativisation}.
    
%     Finally let us consider the fixed point axioms.
%     First notice that a $(\preaxiom)$ axiom,
%     \[
%     \forall y (\phi(\mu \phi,y) \limp y \in \mu \phi)
%     \]
%     has $\pNat$-relativisation,
%     \[
%     \forall y \in \pNat ( \nattrans \phi (\nattrans{(\mu\phi)},y) \limp y \in \nattrans{(\mu\phi)})
%     \]
%     which is logically equivalent to the $(\preaxiom)$ for $\nattrans {(\mu\phi)}$.
%     Next, for an $(\indaxiom)$ axiom,
%     \begin{equation}
%     \label{eq:indaxiom-for-relativisation}
%         \forall x (\phi(\psi,x) \limp \psi(x)) \limp \forall x (x \in \mu\phi \limp \psi(x))
%     \end{equation}
%     we derive its $\pNat$-relativisation as follows,
%     \[
%     \begin{array}{rll}
%          &\forall x ((\pnat x \land \nattrans \phi (\nattrans \psi,x) ) \limp \nattrans \psi(x)) \limp \forall y (y \in \nattrans{(\mu\phi)} \limp \nattrans \psi (y)) & \text{by $(\indaxiom)$ for $\nattrans{(\mu\phi)}$}   \\
%          \implies & \forall x \in \pNat ( \nattrans \phi (\nattrans \psi,x)  \limp \nattrans \psi(x)) \limp \forall y \in \pNat (y \in \nattrans{(\mu\phi)} \limp \nattrans \psi (y)) & \text{by pure logic}
%     \end{array}
%     \]
%     where the last line is just the $\pNat$-relativisation of \eqref{eq:indaxiom-for-relativisation}.
% \end{proof}


\subsection{Relativisation to $\Nat$}
It will be convenient for our realisability argument to work with a notion of \emph{abstract realisability}, where realisability commutes with quantifiers in favour of explicit relativisation of quantifiers to suitable domains \cite{Troelstra98:realizability-handbook}.
The reason for this is that, a priori, all quantifiers of arithmetic are relativised to $\Nat$, but construing so for the fixed point axioms leads to type mismatch during realisability.
For instance, one would naturally like to realise induction axiom for natural numbers by $\Niter$, for which it is natural to consider the $\forall$ of the inductive step unrelativised, whereas the $\forall$ of the conclusion should certainly be relativised to natural numbers.
The same phenomenon presents for the $\indaxiom$ axioms more generally. 
Thus we shall include such relativisations explicitly to handle this distinction.

We introduce a new unary predicate symbol $\pNat $ and write $\langmuarith^\pNat \dfn \langmuarith \cup \{\pNat\}$.\footnote{We use $\pNat$ to avoid confusion with our earlier type $\nat $ for natural numbers.}
$\pNat$ will morally stand for the fixed point $ \mu X\lambda x (x=0 \lor \exists y (Xy \land x=\succ y))$, computing the natural numbers.
However we shall rather take logically equivalent formulations of the prefix and induction axioms for $\pNat$ that are \emph{negative}, in fact ultimately realised by our analogous specialisations of the $\iter{}$ and $\inj{}$ rules for $\nat$ earlier:
\begin{itemize}
    \item $\preaxiom_\pNat^0$: $\pnat \numeral 0$
    \item $\preaxiom_\pNat^\succ$: $\forall x (\pnat x \limp \pnat{ \succ x})$
    \item $\indaxiom_{\pNat,\phi}$:
    $\phi(0) \limp \forall x (\phi(x) \limp \phi(\succ x )) \limp \forall x (\pnat x \limp \phi(x))$
\end{itemize}

We sometimes write $\forall x^\pNat \phi \dfn \forall x (\pnat x \limp \phi)$ indicating that the universal quantifier is \emph{relativised} to $\pNat$.
We similarly write $\exists x^{\pNat} \phi \dfn \exists x (\pnat x \land \phi)$.
In arithmetic all quantifiers are implicitly relativised in this way, by virtue of the induction axiom schema.
However note that the `correct' formulation of induction above, induced by the fixed point definition of $\pNat$, has a non-relativised universal quantifier for the step case. 
This will turn out to be an important refinement in order to avoid type mismatches when conducting realisability. 
For this reason, our realisability argument must work from a theory in which this distinction is native:

\begin{definition}
    Write $\muHAneg$ for the theory defined like $\muHA(\langmuarith^\pNat)$ but, instead of the numerical induction axioms, we include the axioms $\preaxiom_\pNat^0$, $\preaxiom_\pNat^\succ$ and $\indaxiom_\pNat$ defined earlier.
    I.e.\ $\muHAneg$ is the intuitionistic theory over $\langmuarith^\pNat$ including all the axioms of Robinson Arithmetic, all the defining equations of primitive recursive function symbols, all the $\preaxiom_\phi$ and $\indaxiom_{\phi,\psi}$ axioms,  and all the axioms $\preaxiom_\pNat^0$, $\preaxiom_\pNat^\succ$ and $\indaxiom_{\pNat,\phi}$.
\end{definition}


The notation here is suggestive of our analogous notation for negative fragments of $\muLJ$ and $\cmuLJ$ earlier.
Indeed we shall soon see that our notion of `realising type' for $\muHAneg$ will have image over $\nat, \times, \limp,\mu$.
First, let us verify that $\muHAneg$ indeed interprets $\muHA$.


\begin{definition}
    For formulas $\phi$ of $\langmuarith$ we define $\nattrans \phi$ a formula of $\langmuarith^\pNat$ by:
    \begin{itemize}
        \item $ \nattrans{(s=t)} \dfn s=t $
        % \item $\nattrans{(\phi \lor \psi)} \dfn \exists x ((x=0 \limp \nattrans \phi) \land (x = 1 \limp \nattrans \psi))$
        \item $\nattrans {(\phi \star \psi)} \dfn \nattrans \phi \star \nattrans \psi $ for $\star \in\{ \land, \limp\}$
        \item $\nattrans {(\exists x \phi)} \dfn \exists x (\pnat x \land \phi)$
        \item $\nattrans{(\forall x \phi)}\dfn \forall x (\pnat x \limp \phi)$
          \item $ \nattrans{(t \in\mu X \lambda x \phi)} \dfn t \in \mu X \lambda x (\pnat x \land \nattrans\phi)$
    \end{itemize}
\end{definition}


We extend this translation to (definable) predicates, e.g.\ writing $\nattrans \phi(\vec X,\vec x)\dfn \nattrans{\phi(\vec X,\vec x)}$ and, in particular, $\nattrans{(\mu\phi)}x\dfn \nattrans{(x \in \mu\phi)}$.
% We shall employ standard abbreviates for relativised quantifiers, in particular writing $\exists x \in \pNat \, \phi \dfn \exists x (\pnat x \land \phi)$ and $\forall x \in \pNat \, \phi \dfn \forall x (\pnat x \limp \phi) $. 
% 
% 
% Under this interpretation $\muHA$ can be embedded into its induction-free part, a modest extension of pure intuitionistic logic with fixed points. Let us make this clear.
% 
% Write $\muIFOL$ for the intuitionistic theory given by the basic Robinson axioms of arithmetic \todonew{for simplicity. can also just add all true equations} and the $(\preaxiom)$ and $(\indaxiom)$ axioms from \cref{def:muPA}.
% \todonew{can write, say, $\muHA^-$ instead of $\muIFOL$?}
Specialising a well-known reduction from second-order arithmetic to pure second-order logic, we have:

\begin{proposition}
\label{prop:muHA-to-muHAneg}
    If $\muHA \proves \phi$ then $\muHAneg \proves \nattrans \phi$.
\end{proposition}
\begin{proof}
    The Robinson axioms are already part of $\muHAneg$ and, since they are universal statements, so too their relativisations.
    The same argument applies for the defining equations of primitive recursive function symbols.
    
    Each number induction axiom,
    \begin{equation}
        \label{eq:n-ind-instance-for-relativisation}
        \psi(\numeral 0) \limp \forall x (\psi(x) \limp \psi(\succ x)) \limp \forall x \psi(x)
    \end{equation}
    has $\pNat$-translation,
    \[
    \nattrans\psi(\numeral 0)
    \limp 
    \forall x^\pNat (\nattrans\psi(x) \limp \nattrans\psi(\succ x))
    \limp 
    \forall x^\pNat (\nattrans\psi(x))
    \]
    which, by the $\preaxiom_\pNat$ axioms and pure logic, is equivalent to:
    \[
    (\pnat {\numeral 0} \land \nattrans\psi(\numeral 0)
    \limp
    \forall x ((\pnat x \land \nattrans\psi(x)) \limp (\pnat{\succ x } \land \nattrans \psi(\succ x)))
    \limp
    \forall x^\pNat (\pnat x \land \nattrans \psi(x))
    \]
    This is just an instance of $\indaxiom_\pNat$ with invariant $\pnat x \land \nattrans \psi(x)$.

    
    Finally let us consider the fixed point axioms.
    First notice that a $\preaxiom$ axiom,
    \[
    \forall y (\phi(\mu \phi,y) \limp y \in \mu \phi)
    \]
    has $\pNat$-relativisation,
    \[
    \forall y \in \pNat ( \nattrans \phi (\nattrans{(\mu\phi)},y) \limp y \in \nattrans{(\mu\phi)})
    \]
    which is logically equivalent to the $\preaxiom$ axiom for $\nattrans {(\mu\phi)}$.
    
    Next, for an $\indaxiom_\phi$ axiom,
    \begin{equation}
    \label{eq:indaxiom-for-relativisation}
        \forall x (\phi(\psi,x) \limp \psi(x)) \limp \forall x (x \in \mu\phi \limp \psi(x))
    \end{equation}
    we derive its $\pNat$-relativisation as follows,
    \[
    \begin{array}{rll}
         &\forall x ((\pnat x \land \nattrans \phi (\nattrans \psi,x) ) \limp \nattrans \psi(x)) \limp \forall y (y \in \nattrans{(\mu\phi)} \limp \nattrans \psi (y)) & \text{by $\indaxiom$ for $\nattrans{(\mu\phi)}$}   \\
         \implies & \forall x \in \pNat ( \nattrans \phi (\nattrans \psi,x)  \limp \nattrans \psi(x)) \limp \forall y \in \pNat (y \in \nattrans{(\mu\phi)} \limp \nattrans \psi (y)) & \text{by pure logic}
    \end{array}
    \]
    where the last line is just the $\pNat$-relativisation of \eqref{eq:indaxiom-for-relativisation}.
\end{proof}


\subsection{An abstract realisability judgement}
In what follows we shall work with the untyped calculus $\clomuLJneg$ and its typed version $\typedclomuLJneg$ from Section~\ref{sec:term-calculi}.
For convenience we shall employ the following convention:
\begin{convention}
We henceforth identify terms of $\langarith$ with their corresponding definitions in $\muLJneg$. 
We shall furthermore simply identify closed terms of arithmetic with the numerals they reduce to under $\converts$.
Note that, for $s,t$ terms of $\langarith$, we indeed have $s\converts t \implies \muHAneg \proves s=t$.
\end{convention}
At least one benefit of the convention above
is that we avoid any confusion arising from metavariable clash between terms of arithmetic and terms of $\clomuLJneg$. 
This convention is motivated by the clauses of our realisability judgement for atomic formulas.


\medskip



A \emph{(realisability) candidate} is an (infix) relation $\cdot A \cdot $ relating a (untyped) term to a natural number such that $\cdot A n$ is always closed under $\converts$. 
Let us expand $\langmuarith^\pNat $ by construing each realisability candidate as a unary predicate symbol.
The idea is that each $\cdot An $ consists of the just the realisers of the sentence $A\numeral n$.



\anupam{could formulate all of below with closed terms, by defining a valuation function. Need different metavariables, e.g.\ $u,v,w$, reserving $r,s,t$ for terms of $\clomuLJneg$. Do this later if there is time.}
\begin{definition}
    [Realisability judgement]
    For each term $t\in \clomuLJneg$ and closed formula $\phi$ we define the (meta-level) judgement $t \, \realises\, \phi$ as follows:
\begin{itemize}
    % \item $t\, \realises\,  \numeral m = \numeral n$ if $m=n$.\oldtodo{could set to be $\numeral m = t =\numeral n$}
    \item $t\, \realises \, (\numeral m = \numeral n)$ if $\numeral m \converts t \converts \numeral n$.
    \item $t \, \realises \, \pnat {\numeral n}$ if $t\converts \numeral n$.
     \item $t\, \realises\, A\numeral n$ if $t A n$.
    \item $t\, \realises\, \phi_0 \land \phi_1$ if $\projl t \, \realises \phi_0$ and $\projr t \, \realises\, \phi_1$.
    \item $t \, \realises\,  \phi \limp \psi $ if, whenever $s\, \realises \, \phi$, we have $ts \, \realises\, \psi$.
    \item $t\, \realises\, \exists x \phi(x)$ if for some $n\in \Nat$ we have $t\, \realises \, \phi(\numeral n)$.
    \item $t\, \realises\,  \forall x \phi(x)$ if for all $n \in \Nat $ we have $t\, \realises\, \phi(\numeral n)$.
    \item $t\, \realises\, \numeral n \hspace{-.15em}\in \hspace{-.2em} \mu X \lambda x \phi(X,x)$ if $t\, \realises \, \forall x (\phi(A,x)\limp Ax) \limp A\numeral n$ for all candidates $A$.
\end{itemize}
\end{definition}


% \todonew{clean up this remark later}
% \begin{remark}
%     [More about abstract realisability]
    
%     \begin{itemize}
%         \item to avoid type mismatches, and to expose our moral specialisation of HA2 to F.
%         \item note that Berger and Tsuiki employ a similar approach, albeit with a different realisability model interpreting fixed points as fixed points, akin to our earlier type structure $\hr\cdot$.
%         \item note also that the choice of abstract realisability is consistent, morally, with our treatment of $\mu$ formulas under a second-order encoding, where we simply commute with the universal quantification over candidates under type erasure.
%     \end{itemize}
% \end{remark}


\begin{definition}
    [Realising type]
    The \emph{realising type} of a (possibly open) formula $\phi$ without candidate symbols, written $\type \phi$, is given by:
\begin{itemize}
    \item $\type{s=t} \dfn \N$
    \item $\type{\pnat t} \dfn \N$
    \item $\type {Xt} \dfn X$
    \item $\type{\phi\land \psi} \dfn \type \phi \times \type \psi$
    \item $\type{\phi \limp \psi} \dfn \type \phi \to \type \psi$
    \item $\type{\exists x \phi}\dfn  \type \phi$
     \item $\type{\forall x \phi} \dfn \type\phi$
    \item $\type {t \in \mufo X x \phi} \dfn \mu X  \type \phi$
    % \item $\type{t \in \mu X \lambda x \phi} \dfn \muN X\, \type \phi $
\end{itemize}
\end{definition}




\subsection{Closure under $\converts$ and realising induction}
Our realisability model is compatible with our notion of conversion from $\muLJ$:

\begin{lemma}
    [Closure under $\converts$]
    \label{lem:realises-closed-under-converts}
    If $t \, \realises\, \phi$ and $t\converts t'$ then $t'\, \realises \, \phi$.
\end{lemma}
\begin{proof}
    By induction on the structure of $\phi$:
    \begin{itemize}
        \item if $\phi$ is $\numeral m=\numeral n$ then still $\numeral m \converts t' \converts \numeral n$ by symmetry and transitivity of $\converts$.
        \item if $\phi$ is $\pnat {\numeral n}$ then still $t' \converts \numeral n$ by symmetry and transitivity of $\converts$.
        \item if $\phi$ is $A\numeral n $ then still $t' A n$ by definition of candidate.
         \item if $\phi$ is $\phi_0 \land \phi_1$ then we have $\projl t \, \realises\, \phi_0$ and $\projr t\, \realises\, \phi_1$, and so by IH and context-closure of $\converts$ we have $\projl t' \, \realises\, \phi_0$ and $\projr t'\, \realises\, \phi_1$, thus indeed $t'\, \realises\, \phi$.
        \item if $\phi$ is $\phi_0 \limp \phi_1$ then for any $s\, \realises\, \phi_0$ we have $ts \, \realises \, \phi_1$, and so also $t's \, \realises \, \phi_1$ by IH and context-closure of $\converts$. 
        Since choice of $s$ was arbitrary we are done.
        \item if $\phi $ is $\exists x\phi'(x)$ then there is some $n\in \Nat $ with $ t \, \realises\, \phi'(\numeral n)$. 
        So by IH we have $ t' \, \realises\, \phi'(\numeral n)$, and so also $t'\, \realises\, \phi $.
        \item if $\phi$ is $\forall x \phi'(x)$ then for each $n\in \Nat$ we have $t\, \realises\, \phi'(\numeral n)$, and so also $t' \, \realises\, \phi'(\numeral n)$ by IH. Since choice of $n$ was arbitrary we are done.
        \item if $\phi$ is $\mu\phi'\, \numeral n$ then, for all candidates $A$ we have $t \, \realises \, \forall x (\phi'(A,x) \limp Ax) \limp A\numeral n$, and so also $t' \, \realises \, \forall x (\phi'(A,x) \limp Ax) \limp A\numeral n$ by IH.
        Since choice of $A$ was arbitrary we are done. \qedhere
    \end{itemize}
\end{proof}

Note that, despite its simplicity, the above lemma has the following consequence, by consideration of the candidate $\{(t,n) : t \, \realises\, \psi(n)\}$:
\begin{lemma}
\label{prop:mu-realiser-realises-formula-invariants}
    If $t \, \realises\, \mu \phi\, \numeral n$ then, for any formula $\psi(x)$, $t \, \realises\, \forall x (\phi(\psi,x) \limp \psi(x)) \limp \psi(\numeral n)$.
\end{lemma}

This allows us to realise all the induction axioms for fixed points:
\begin{proposition}
    [Realising Induction]
    \label{prop:iter-realises-ind}
    Let $\sigma(X) = \type {\phi(X,x)}$ and $\tau = \type{\psi(x)}$. 
    Then  $\iter{\sigma,\tau}\, \realises\, \indaxiom_{\phi,\psi}$.
\end{proposition}
\begin{proof}
    We shall omit subscripts to lighten the syntax.
    Let $u\, \realises\, \forall x (\phi(\psi,x) \limp \psi(x))$ and $v \, \realises\, \mu\phi\, \numeral n$.
    We need to show that $\iter{}\, u\, v\, \realises\, \psi(n)$:
    \[
    \begin{array}[b]{r@{\ \realises\ }ll}
         v & \forall x (\phi(\psi,x) \limp \psi(x)) \limp \psi(\numeral n) & \text{since $v\, \realises\, \mu\phi\, \numeral n$ and by \Cref{prop:mu-realiser-realises-formula-invariants}} \\
         v\, u & \psi(\numeral n) & \text{by $\realises\limp$ since $u\, \realises \, \forall x (\phi(\psi,x) \limp \psi(x))$}\\
         \iter\, u\, v & \psi(\numeral n) & \text{by $\converts$ and \Cref{lem:realises-closed-under-converts}}  \tag*{\qed}
    \end{array}
    \]
\renewcommand{\qed}{}
\end{proof}


\subsection{Functoriality and realising prefix axioms}

As expected we will realise $\preaxiom$ by $\inj$.
For this we need a `functoriality lemma' establishing a semantics for functors of $\muLJ$ within our realisability model.
This is because the reductions for $\inj$ introduce functors.

In fact, due to the fact that we admit inductive types as primitive, with native $\lr\mu$ and $\rr \mu$ rules, we shall have to establish the realisability of $\preaxiom$ and the functoriality properties \emph{simultaneously}.


\begin{lemma}
    [Functoriality and realising $\preaxiom$]
    \label{lem:funct-and-pre-realisation}
    Let $\type{\phi(\vec Y,Z,z)} = \sigma(\vec Y,Z)$.
    \todonew{G: Write simply $\mu\phi(\vec Y)$ for $\mu Z\lambda z \phi(\vec Y,Z,z)$}
    Then for all $\vec \tau = \type{\vec \psi(x)}$ we have:
    \begin{enumerate}
        \item\label{item:in-realises-pre} $\injX{\sigma(\vec \tau)} \, \realises\, \preaxiom_{\phi(\vec \psi)} $
        \item\label{item:funct-lem} Let $t \, \realises\, \forall z (\chi(z) \limp \chi'(z))$ \todonew{G: with $\type{\chi(z)} = \gamma$ and $\type{\chi'(z)} = \gamma'$}. Then:
        \todonew{G: THE RED TEXT CAN BE MOVED IN THE PROOF}
        \begin{enumerate}
            \item\label{item:funct-lem-pos} if $\phi(\vec Y,Z,z)$ is positive in $Z$ then $\sigma(\vec \tau,t)\, \realises\, \forall z (\phi(\vec \psi,\chi,z) \limp \phi (\vec\psi , \chi',z))$.
            \item\label{item:funct-lem-neg}
            if $\phi(\vec Y,Z,z)$ is negative in $Z$ then $\sigma(\vec \tau,t)\, \realises\, \forall z (\phi(\vec \psi,\chi',z) \limp \phi (\vec\psi , \chi,z))$.
        \end{enumerate}
    \end{enumerate}
\end{lemma}

Before proving this, let us make some pertinent remarks:
\begin{remark}
    [Comparison to $\F$]
    For the reader used to fixed points via second-order encodings, it is perhaps surprising that functoriality and realisability of $\preaxiom$ are mutually dependent. 
    To recall, in system $\F$, we usually set $\mu X\sigma(X) \dfn \forall X ((\sigma(X) \to X) \to X)$. 
    From here we might set the appropriate functor as,
    \[
        \mu X \sigma(X,t) \dfn \Lambda X ((\sigma(X,t)\to X)\to X)
    \]
    whence the appropriate functoriality properties, analogous to \eqref{item:funct-lem} of \Cref{lem:funct-and-pre-realisation}, are readily established directly, without appeal to the realisation of $\preaxiom$.
    However the above is not (explicitly) how $\muLJ$ defines the fixed point functors.
    
    While a second-order calculus introduces $\forall$ by way of an eigenvariable representing an \emph{arbitrary} prefixed point, our $\rr\mu$ rule works precisely because $\mu X\sigma(X)$ is the \emph{least} fixed point, and so is already sufficiently general.
    This can be realised by appropriate proof transformations, themselves relying on the $\rr\mu$ rule.
    It is for this reason that the mutual dependency above presents in our setting.
\end{remark}

\begin{proof}
[Proof of \Cref{lem:funct-and-pre-realisation}]
    We proceed by induction on the structure of $\phi(\vec Y,Z,z)$, proving both \eqref{item:in-realises-pre} and \eqref{item:funct-lem} simultaneously.
    More precisely, \eqref{item:in-realises-pre} will rely on instances of \eqref{item:funct-lem} for the same $\phi$, and \eqref{item:funct-lem} will rely on smaller instances of both \eqref{item:in-realises-pre} and \eqref{item:funct-lem}.


    % Let us simply write $\mu\phi(\vec \psi)$ for $\mu Z\lambda z \phi(\vec \psi,Z,z)$, so that 
    We have $\preaxiom_{\phi(\vec \psi)} = \forall z (\phi(\vec \psi,\mu\phi(\vec \psi),z) \limp \mu\phi(\vec \psi)\, z)$.
    So to prove \eqref{item:in-realises-pre} let, 
    \begin{equation}
        \label{eq:in-realises-pre-lfp-unfolded-realiser}
        u\, \realises \, \phi(\vec \psi,\mu\phi(\vec \psi),\numeral n)
    \end{equation}
    and, by $\realises\forall $ and $\realises\limp$, we need to show $\injX{\sigma(\vec \tau)}\, u\, \realises\, \mu\phi(\vec \psi)\, \numeral n$.
    Now, by $\realises\mu$ and $\realises\limp$, let $C$ be a candidate and, 
    \begin{equation}
    \label{eq:in-realises-pre-candidate-prefix-realiser}
        v\, \realises \, \forall z (\phi(\vec \psi,C,z) \limp Cz)
    \end{equation} 
    so that it suffices to show $ \injX{\sigma(\vec \tau)}\, u\, v\, \realises\, C\numeral n$.
    We have:\footnote{The types for $\iter{}$ are omitted but determined by context.}
    \[
    \begin{array}{r@{\ \realises \ }ll}
         \iter{}\, v & \forall z (\mu\phi(\vec \psi)\, z \limp Cz) & \text{by \Cref{prop:iter-realises-ind}, \eqref{eq:in-realises-pre-candidate-prefix-realiser} and $\realises\limp$} \\
         \sigma (\vec \tau,\iter{} \, v ) & \phi(\vec \psi,\mu\phi(\vec \psi),\numeral n) \limp \phi(\vec \psi,C,\numeral n) & \text{by IH for \cref{item:funct-lem} and $\realises \forall$} \\
         \sigma (\vec \tau,\iter{} \, v ) \, u & \phi(\vec \psi,C,\numeral n) & \text{by \eqref{eq:in-realises-pre-lfp-unfolded-realiser} and $\realises \limp$} \\
         v (\sigma (\vec \tau,\iter{} \, v ) \, u) & C\numeral n & \text{by \eqref{eq:in-realises-pre-candidate-prefix-realiser} and $\realises \forall $ and $\realises\limp$} \\
         \injX{\sigma(\vec \tau)}\, u\, v & C\numeral n & \text{by $\converts \inj$ and \Cref{lem:realises-closed-under-converts}}
    \end{array}
    \]

    To prove \eqref{item:funct-lem} the critical case is when $\phi(\vec Y,Z,z) $ is a fixed point formula $ \numeral m \in \mu X \lambda x \phi'(\vec Y,\allowbreak Z,z,X,x)$.
    We consider only the case when $Z$ is positive, \eqref{item:funct-lem-pos}.
    As before we shall simply write, say $\mu\phi'(\vec Y,Z,z)$ for $\mu X\lambda x \phi'(\vec \psi, Z,z,X,x)$.
    To prove \eqref{item:funct-lem-pos} let $n \in \Nat$ and
    \begin{equation}
    \label{eq:funct-lem-mu-antecedent}
        u\ \realises\ \phi(\vec \psi,\chi,\numeral n)
        \qquad\qquad
        \text{(i.e.\ $u\, \realises\, \mu\phi'(\vec \psi,\chi,\numeral n)\, \numeral m$)}
    \end{equation}
    so, by $\realises\forall$ and $\realises\limp$ we need to show $\sigma(\vec \tau,t)\, u\,  \realises\, \phi(\vec \psi,\chi',\numeral n)$, i.e.\ $ \sigma(\vec \tau,t)\, u\,  \realises\,  \mu\phi'(\vec \psi,\chi',\numeral n)\, \numeral m$.
    So, by $\realises\mu$, let $A$ be a candidate and:
    \begin{equation}
    \label{eq:funct-lem-mu-succedent-input}
        v\ \realises\ \forall x (\phi'(\vec \psi,\chi',\numeral n, A,x) \limp Ax)
    \end{equation}
    We need to show that $\sigma(\vec \tau,t)\, u\, v\, \realises\, A\numeral m$.


    Now let $\sigma'(\vec Y,Z,X) = \type {\phi'(\vec Y,Z,z,X,x)}$. 
    Again we may similarly write simply $\mu\sigma'(\vec Y,Z)$ for $\mu X \sigma'(\vec Y,Z,X)$ (which is just $\sigma(\vec Y,Z)$).
    First, by IH\eqref{item:in-realises-pre} for $\sigma'(\vec Y,Z,X)$, we have:
    \begin{equation}
    \label{eq:in-realises-pre-ih}
        \injX{\sigma'(\vec \tau,\gamma')} \, \realises\, \forall x (\phi'(\vec \psi, \chi',\numeral n, \mu\phi'(\vec \psi,\chi',\numeral n),x) \limp \mu\phi'(\vec \psi,\chi',\numeral n)\, x)
    \end{equation}
    We claim that:
    \begin{equation}
    \label{eq:mur-realises-pre-with-context}
    \rr\mu\, \sigma'(\vec \tau,t,\sigma(\vec \tau,\gamma'))\ \realises\ 
    \forall x (\phi'(\vec \psi,\chi,\numeral n,\mu\phi'(\vec \psi,\chi',\numeral n),x) \limp \mu\phi'(\vec \psi,\chi',\numeral n)\, x)
    \end{equation}
    To prove this let $w\, \realises\, \phi'(\vec \psi, \chi,\numeral n,\mu\phi'(\vec \psi,\chi',\numeral n),\numeral k)$ and we show $ \rr\mu\, \sigma'(\vec \tau,t,\sigma(\vec \tau,\gamma'))\, w \, \realises\, \mu\phi'(\vec \psi,\allowbreak \chi',\numeral n)\, \numeral k$:
    \[
    \begin{array}{r@{\ \realises\ }ll}
         % \sigma'(\vec \tau,t,\mu X \sigma(\vec \tau,\gamma')) & \forall z (\phi'(\vec \psi,\chi,z, \mu\phi'(\vec \psi,\chi',\numeral n), \numeral k) \limp \phi'(\vec \psi, \chi', z, \mu\phi'(\vec \psi,\chi',\numeral n) , \numeral k))
         % & \text{by IH\eqref{item:funct-lem} for $\sigma'(\vec Y,Z,X)$}
         % \\
         \sigma'(\vec \tau,t,\mu X \sigma(\vec \tau,\gamma')) \, w & \phi'(\vec \psi,\chi',\numeral n, \mu\phi'(\vec \psi, \chi',\numeral n), \numeral k) &\text{by IH\eqref{item:funct-lem} for $\phi'$, $\realises\forall $ and $\realises\limp$} \\
         \injX{\sigma'(\vec \tau,\gamma')} (\sigma'(\vec \tau,t,\mu X \sigma(\vec \tau,\gamma')) \, w) & \mu\phi'(\vec \psi,\chi',\numeral n)\, \numeral k & \text{by \eqref{eq:in-realises-pre-ih} and $\realises\forall$ and $\realises\limp$} \\
         \mu_r\, \sigma'(\vec \tau, t, \sigma(\vec \tau,\gamma'))\, w & \mu\phi'(\vec \psi, \chi', \numeral n)\, \numeral k & \text{by $\converts$}
    \end{array}
    \]
    Now by \eqref{eq:mur-realises-pre-with-context}, $\realises\mu$ and \Cref{prop:mu-realiser-realises-formula-invariants} we have,
    \begin{equation}
        u\, \realises\, \forall x (\phi'(\vec \psi,\chi,\numeral n,\mu\phi'(\vec \psi,\chi',\numeral n),x) \limp \mu\phi'(\vec \psi,\chi',\numeral n)\, x) \limp \mu\phi'(\vec \psi,\chi',\numeral n)\, \numeral m
    \end{equation}

    Now we prove $\sigma(\vec \tau,t)\, u\, v\, \realises\, A\numeral m$ as follows:
    \[
    \begin{array}{r@{\ \realises \ }ll}
          % \rr\mu\, \sigma'(\vec \tau,t,\sigma(\vec \tau,\gamma')) & \forall x (\phi'(\vec \psi,\chi,\numeral n,\mu\phi'(\vec \psi,\chi',\numeral n),x) \limp \mu\phi'(\vec \psi,\chi',\numeral n)\, x) & \text{by \eqref{eq:mur-realises-pre-with-context}} \\
          u\, (\rr\mu\, \sigma'(\vec \tau,t,\sigma(\vec \tau,\gamma'))) & \mu\phi'(\vec \psi,\chi',\numeral n)\, \numeral m & \text{by \eqref{eq:mur-realises-pre-with-context} and $\realises\limp$}\\
          \iter{}\, (\rr\mu\, \sigma'(\vec \tau,t,\sigma(\vec \tau,\gamma')))\, u & \mu\phi'(\vec \psi,\chi',\numeral n)\, \numeral m & \text{by $\converts$} \\
          \iter{}\, (\rr\mu\, \sigma'(\vec \tau,t,\sigma(\vec \tau,\gamma')))\, u\, v & A \numeral m & \text{by $\realises \mu$ and \eqref{eq:funct-lem-mu-succedent-input}} \\
          \mu \sigma'(\vec \tau,t) \, u \, v & A\numeral m & \text{by defintion of $\mu$-functor} \\
          \sigma(\vec \tau,t)\, u\, v & A\numeral m
    \end{array}
    \]


    The remaining steps for functoriality, \eqref{item:funct-lem}, are routine, and for these we shall simply suppress the parameters $\vec \psi,\vec \tau$ henceforth, writing simply $\phi(Z,z)$ instead of $\phi(\vec Y,Z,z)$ and $\sigma(Z)$ instead of $\sigma(\vec Y,Z)$.
    This is because we shall never vary the parameters $\vec \psi,\vec \tau$ in the remaining cases.
    We give only the cases for \cref{item:funct-lem-pos}, the ones for \cref{item:funct-lem-neg} being similar.
     We need to show that $\sigma(t) \, \realises\, \forall z (\phi(\chi,z) \limp \phi(\chi',z))$ so let $n\in \Nat$ and $u\, \realises\, \phi(\chi,\numeral n)$ and let us show that $\sigma(t)\, u\,  \realises\, \phi(\chi',\numeral n)$, under $\realises\forall$ and $\realises\limp$.

    
    \begin{itemize}
        \item If $\phi(Z,z)=Zs(z)$ then $\sigma(t) = t $. So we need $t\, \realises\, \forall z(\chi(s(z)) \limp \chi'(s(z))) $. Let $n\in \Nat$ and, since already $t\, \realises\, \forall z (\chi(z) \limp \chi'(z))$ by assumption, we have $t\, \realises \chi(\numeral{s(n)}) \limp \chi(\numeral {s(n)})$ as required.
        \item If $\phi(Z,z)$ is any other atomic formula $\theta$ then $\forall z (\phi(\chi,z)\limp \phi(\chi',z))$ is just $\forall z (\theta \limp \theta)$ and $\sigma(t)$ is just $\id\, \realises \, \forall z (\theta \limp \theta)$ by $\realises\forall$ and $\id\converts$.
        \item Suppose $\phi(Z,z)$ is $\phi_0(Z,z) \land \phi_1(Z,z)$. 
        Write $\sigma_i(Z) \dfn \type {\phi_i(Z,z)}$ for $i\in \{0,1\}$ so that $\sigma(Z) = \sigma_0(Z) \times \sigma_1(Z)$.
        % We have by IH that:
        % \begin{equation}
        % \label{eq:funct-lem-ih-conj}
        %     \sigma_i(t)\, \realises\, \forall z (\phi_i (\chi,z)\limp \phi_i (\chi',z))
        % \end{equation}
        % for $i\in \{0,1\}$.
        We need to show that $\sigma(t)\, \realises\, \forall z (\phi(\chi,z) \limp \phi(\chi',z))$, so let $n\in \Nat$ and $u\, \realises\, \phi(\chi,\numeral n)$ and let us show that $\sigma(t)\, u \, \realises \, \phi(\chi',\numeral n)$:
        \[
        \begin{array}{r@{\ \realises\ }ll}
             \proj i u & \phi_i(\chi,\numeral n) & \text{for $i\in \{0,1\}$ by $\realises\land$} \\
             \sigma_i(t) (\proj i u) & \phi_i(\chi', \numeral n) & \text{for $i\in \{0,1\}$ by IH and $\realises\forall$ and $\realises\limp$} \\
             \proj i (\sigma(t)\, u) & \phi_i(\chi',\numeral n) &\text{for $i\in \{0,1\}$ by form of $\sigma(t)$ and $\converts$} \\
             \sigma(t)\, u  & \phi(\chi',\numeral n) & \text{by $\realises\land$}
        \end{array}
        \]
        \item Suppose $\phi(Z,z)$ is $\phi_0(Z,z) \limp \phi_1(Z,z)$ with $\phi_0(Z,z)$ negative in $Z$ and $\phi_1(Z,z)$ positive in $Z$. 
        Write $\sigma_i(Z)\dfn \type{\phi_i(Z,z)}$ for $i\in \{0,1\}$ so that $\sigma(Z) = \sigma_0(Z) \limp \sigma_1(Z)$.
        We need to show that $\sigma(t)\, \realises\, \forall z (\phi(\chi,z) \limp \phi(\chi',z))$ so let $n\in \Nat $ and $u\, \realises\, \phi(\chi,\numeral n)$ and let us show that $\sigma(t)\, u \, \realises\, \phi(\chi',\numeral n)$.
        By the form of $\phi(Z,z)$, and by $\realises\limp$, let $v\, \realises\, \phi_0(\chi',\numeral n)$ so that it suffices to show $\sigma(t)\, u\, v\, \realises\, \phi_1(\chi', \numeral n)$:
        \[
        \begin{array}{r@{\ \realises\ }ll}
             \sigma_0(t)\, v & \phi_0 (\chi,\numeral n) & \text{by IH and $\realises\forall $ and $\realises\limp$} \\
             u\, (\sigma_0(t)\, v) & \phi_1(\chi,\numeral n) & \text{by $\realises\limp$} \\
             \sigma_1(t)\, (u\, (\sigma_0(t)\, v)) & \phi_1(\chi',\numeral n) & \text{by IH and $\realises\forall$ and $\realises\limp$} \\
             \sigma(t)\, u\, v & \phi_1(\chi',\numeral n) & \text{by form of $\sigma(t)$ and $\converts$}
        \end{array}
        \]
        
        \item Suppose $\phi(Z,z) $ is $\exists y \phi'(Z,z,y)$. Write $\sigma'(Z)\dfn \type{\phi'(Z,z,y)}$ so that $\sigma(Z) = \sigma'(Z)$.
        We need to show that $\sigma(t) \, \realises\, \forall z (\phi(\chi,z) \limp \phi(\chi',z))$ so let $n\in \Nat$ and $u\, \realises\, \phi(\chi,\numeral n)$ and let us show that $\sigma(t)\, u\,  \realises\, \phi(\chi',\numeral n)$, under $\realises\forall$ and $\realises\limp$. 
        \[
        \begin{array}{r@{\ \realises\ }ll}
             u & \exists y \phi'(\chi,\numeral n,y) & \text{by assumption and form of $\phi(Z,z)$} \\
             u & \phi' (\chi,\numeral n, \numeral k) & \text{for some $k\in \Nat$ by $\realises\exists$} \\
             \sigma'(t)\, u & \phi'(\chi',\numeral n, \numeral k) & \text{for some $k\in \Nat$ by IH and $\realises\forall$ and $\realises\limp$} \\
             \sigma'(t)\, u & \exists y \phi'(\chi', \numeral n, y) & \text{by $\realises\exists$} \\
             \sigma(t)\, u & \phi(\chi',\numeral n) & \text{by forms of $\sigma(Z)$ and $\phi(Z,z)$}             
        \end{array}
        \]
        \item Suppose $\phi(Z,z) $ is $\forall y \phi'(Z,z,y)$. Write $\sigma'(Z) \dfn \type{\phi'(Z,z,y)}$ so that $\sigma(Z) = \sigma'(Z)$.
        We need to show that $\sigma(t)\, \realises\, \forall z (\phi(\chi,z) \limp \phi(\chi',z))$ so let $n\in \Nat $ and $u\, \realises\, \phi(\chi,\numeral n)$ and let us show that $\sigma(t)\, u\, \realises\, \phi(\chi',\numeral n)$ under $\realises\forall$ and $\realises\limp$.
        \[
        \begin{array}[b]{r@{\ \realises\ }ll}
             u & \forall y \phi'(\chi,\numeral n,y) & \text{by assumption and form of $\phi(Z,z)$} \\
             u & \phi'(\chi,\numeral n, \numeral k) & \text{for all $k\in \Nat$ by $\realises \forall$} \\
             \sigma'(t)\, u & \phi'(\chi',\numeral n, \numeral k) & \text{for all $k\in \Nat$ by IH and $\realises\forall$ and $\realises\limp$} \\
             \sigma'(t)\, u & \forall y \phi'(\chi',\numeral n, y) & \text{by $\realises\forall$} \\
             \sigma(t)\, u & \phi(\chi', \numeral n) & \text{by forms of $\sigma(Z)$ and $\phi(Z,z)$} \tag*{\qed}
        \end{array}
        \]
    \end{itemize}
    \renewcommand{\qed}{}
\end{proof}




\subsection{Putting it all together}

\begin{theorem}
    [Soundness]
    \label{thm:realisability-soundness}
    If $\muHAneg \proves \phi$ then there is a typed $\typedclomuLJneg$ term $t:\type \phi$ such that $t\, \realises\, \phi$.
\end{theorem}
\begin{proof}
The axioms and rules of intuitionistic predicate logic are realised as usual, not using the fixed point rules. 
The basic arithmetical axioms of $\muHAneg$ are just universal closures of true equations between primitive recursive terms, which are realised by the corresponding derivations of those terms in $\muLJneg$.\footnote{Recall that $\muLJneg$ represents at least all the (even higher-order) primitive recursive functions.}
The axiom schemes $\preaxiom$ and $\indaxiom$ are realised by $\inj$ and $\iter{}$ of appropriate types, by \Cref{lem:funct-and-pre-realisation} and \Cref{prop:iter-realises-ind} respectively.
    It remains to consider the $\pNat$-specific axioms,  which are realised by the analogous $\nat$-specific rules of $\muLJneg$:
    \begin{itemize}
        \item $\numeral 0 \, \realises\, \nprezero$, as this is just $\numeral 0\, \realises\, \pnat {\numeral 0}$, which follows by definition of $\realises\pNat$.
        \item $\succ \, \realises\, \npresucc$. Let $n\in \Nat $ and $u\, \realises\, \pnat {\numeral n}$, i.e.\ $u\converts \numeral n$, so that, by $\realises\forall $ and $\realises\limp$, it suffices to show that $\succ u\, \realises\, \pnat{\succ \numeral n}$.
        We have that $\succ u \converts \numeral{n+1} = \succ \numeral n $, and so indeed $\succ u\, \realises\, \pnat{\succ \numeral n}$ as required.
        \item $\Niter\, \realises \, \nindaxiom$.
        Let $u\, \realises\, \phi(\numeral 0)$ and $v\, \realises\, \forall x (\phi(x)\limp \phi(\succ x ))$ so that it suffices, by $\realises\limp$, to show that $\Niter u\, v\, \realises \, \forall x^\pNat \phi(x)$.
        For this let $n\in \Nat $ and $w\, \realises\, \pnat {\numeral n}$, and we show $\Niter u\, v\, w\, \realises\, \phi(\numeral n)$ by induction on $n$:
        \begin{itemize}
            \item If $n=0$ then $w\converts \numeral 0$ so $\Niter u\, v\, w\, \converts\, \Niter u\, v\, 0\, \converts u\, \realises \, \phi(\numeral 0)$ by assumption.
            \item If $n = n'+1$ then we have:
            \[
            \begin{array}[b]{r@{\ \realises\ }ll}
                 \Niter u\, v\, \numeral n' & \phi(\numeral n') & \text{by IH} \\
                 v\, (\Niter u\, v\, \numeral n') & \phi(\succ{\numeral n'}) & \text{by $\realises\forall$ and $\realises\limp$} \\
                 \Niter u\, v\, w & \phi(\numeral n) & \text{by $\converts$} \tag*{\qed}
            \end{array}
            \]
        \end{itemize}
    \end{itemize}
    \renewcommand{\qed}{}
\end{proof}

\begin{corollary}
    Any provably total recursive function of $\muPA$ is representable in $\muLJneg$.
\end{corollary}
\begin{proof}
    Assume $\muPA \proves \forall \vec x \exists ! y \, s(\vec x, y) = t(\vec x, y)$, without loss of generality.\footnote{Recall we have included each primitive recursive function definition as a symbol of $\langarith$ and defining equations among the axioms of $\PA$, $\HA$ and extensions.} 
    So we have,
    \[
    \begin{array}{rll}
          & \muPA \proves \forall \vec x \exists  y \, s(\vec x,y) = t(\vec x,y) \\
          \implies & \muHA \proves \forall \vec x \exists  y \, s(\vec x,y) = t(\vec x,y) & \text{by \Cref{thm:muPA-pi02-cons-muHA}} \\
           \implies  & \muHAneg \proves \forall \vec x^\pNat \exists y^\pNat s(\vec x,y) = t(\vec x,y) & \text{by \Cref{prop:muHA-to-muHAneg}} \\
         %    \implies  & \muHAneg \proves \forall \vec x^\pNat \exists y\,  s(\vec x,y) = t(\vec x,y) & \text{by pure logic} \\
         % \implies  &  u\, \realises \forall \vec x^\pNat \exists y\,  s(\vec x,y) = t(\vec x,y)  &  \text{for some $ u:\vec \nat \to \nat $ in $ \typedclomuLJneg$ by \Cref{thm:realisability-soundness}} \\
         %  \implies  & u\, \realises \forall \vec x^\pNat \exists y\,  s(\vec x,y) = t(\vec x,y) & \text{for some $u:\vec \nat \to \nat $ in $ \muLJneg$ by \Cref{thm:terms-to-derivations}} \\    \implies & \forall \vec m \in \Nat\, u\, \vec {\numeral m}\, \realises \, \exists y\, s(\vec {\numeral m},y) = t(\vec {\numeral m},y) & \text{by $\realises \forall$, $\realises \limp$ and $\realises\pNat$} \\
         %   \implies & \forall \vec m \in \Nat \exists n \in \Nat\,  u\, \vec {\numeral m}\, \realises \,  s(\vec {\numeral m},\numeral n) = t(\vec {\numeral m},\numeral n) & \text{by $\realises \exists$} \\
          \implies  &  u\, \realises\,  \forall \vec x^\pNat \exists y^\pNat s(\vec x,y) = t(\vec x,y) 
        \end{array}
          \]
          for some $u$ of $ \typedclomuLJneg$ (of appropriate type) by \Cref{thm:realisability-soundness}.
                    Thus:
          % &  \text{for some $ u $ of $ \typedclomuLJneg$ by \Cref{thm:realisability-soundness}} \\
          % \implies  & u\, \realises \forall \vec x^\pNat \exists y^\pNat s(\vec x,y) = t(\vec x,y) & \text{for some $u $ of $ \muLJneg$ by \Cref{thm:terms-to-derivations}} \\
          \[
              \begin{array}{rll}
          & \forall \vec m \in \Nat\ u\, \vec {\numeral m}\, \realises \, \exists y^\pNat s(\vec {\numeral m},y) = t(\vec {\numeral m},y) & \text{by $\realises \forall$, $\realises \limp$ and $\realises\pNat$} \\
          \implies & \forall \vec{ m} \in \Nat \exists n \in \Nat\, (\projl (u\, \vec{\numeral m})\, \realises\, \pnat {\numeral n} \text{ and } \projr (u\, \vec {\numeral m}) \, \realises \, s(\vec {\numeral m}, \numeral n) = t(\vec {\numeral m}, \numeral n) ) & \text{by $\realises \exists$, $\realises\land$ and $\realises\pNat$} \\
          \implies & \forall \vec{ m} \in \Nat \exists n \in \Nat\, (\projl (u\, \vec {\numeral m}) \converts \numeral n \text{ and } s(\vec {\numeral m},\numeral n) \converts t(\vec {\numeral m}, \numeral n)) &\text{by $\realises\pNat$ and $\realises=$}
    \end{array}
    \]
    Thus, as the graph of $s(\vec x,y) = t(\vec x,y)$, in the standard model $\nmod$ is functional, by assumption, it is computed by $ \lambda \vec x^{\vec \nat} (\projl (u\, \vec x))$.
By \Cref{thm:terms-to-derivations} this is equivalent to some $\muLJneg$ derivation too.
\todonew{G: WE COULD REPLACE $t(\vec x, y)$ with $0$.}
\end{proof}